\documentclass[13pt,a4paper]{extarticle}
\usepackage{fontspec}
\usepackage{unicode-math}
\setmainfont{Latin Modern Roman}
\setmathfont{Latin Modern Math}
\usepackage{amsmath}
\usepackage[margin=1.8cm]{geometry}
\usepackage{enumitem}
\usepackage{array}
\usepackage{multicol}
\setlength{\columnsep}{1.2em}
\usepackage{eurosym}
\linespread{1.04}
\renewcommand{\arraystretch}{1.28}
\usepackage{newunicodechar}
\newunicodechar{·}{\kern-.13em\textperiodcentered\kern-.13em}
\usepackage{parskip}
\setlength{\parindent}{0pt}
\begin{document}
\section*{Exercicis i problemes}

\begin{quote}\small Relació d'exercicis del bloc de \textbf{Nombres reals}, ordenada seguint l'ordre dels apunts i amb repàs acumulatiu: a mesura que avança, incorpora continguts dels apartats anteriors. Els marcats \emph{(aplicat)} parteixen d'un context real o d'una altra matèria; els marcats \emph{(repàs)} en combinen diversos.\end{quote}

[Descarrega el llistat en PDF](../assets/exercicis/1\_nombres\_reals.pdf)\par

\vspace{1mm}\hrule\vspace{3mm}

\section*{1. El conjunt dels nombres reals i la recta real}

\par\noindent\begin{minipage}{\linewidth}
\textbf{1.} Classifica cada nombre indicant a quins conjunts pertany ($\mathbb{N}$, $\mathbb{Z}$, $\mathbb{Q}$, $\mathbb{R}\setminus\mathbb{Q}$; un nombre pot pertànyer a més d'un):\par

\begin{tabular}{@{}p{.46\linewidth}p{.46\linewidth}@{}}
\textbf{a)} $-7$ & \textbf{b)} $\dfrac{3}{4}$ \\[2pt]
\textbf{c)} $\sqrt{16}$ & \textbf{d)} $\sqrt{5}$ \\[2pt]
\textbf{e)} $0$ & \textbf{f)} $2{,}5$ \\[2pt]
\textbf{g)} $-\dfrac{9}{3}$ & \textbf{h)} $\pi$ \\[2pt]
\end{tabular}\par
\end{minipage}\par\vspace{1.6mm}
\par\noindent\begin{minipage}{\linewidth}
\textbf{2.} Digues si cada afirmació és certa o falsa, i justifica-ho breument:\par

\begin{itemize}[leftmargin=0pt,itemsep=3pt,topsep=3pt,parsep=0pt,label={}]
\item \textbf{a)} $\mathbb{N}\subset\mathbb{Z}$
\item \textbf{b)} $\sqrt{2}\in\mathbb{Q}$
\item \textbf{c)} Tot nombre enter és racional.
\item \textbf{d)} $0{,}\overline{6}$ és irracional.
\end{itemize}
\end{minipage}\par\vspace{1.6mm}
\par\noindent\begin{minipage}{\linewidth}
\textbf{3.} Ordena de més petit a més gran: $-3$, $\dfrac{7}{2}$, $-1{,}5$, $4$, $0$.\par
\end{minipage}\par\vspace{1.6mm}
\par\noindent\begin{minipage}{\linewidth}
\textbf{4.} Sense calculadora, decideix quin nombre és més gran, $\dfrac{2}{3}$ o $0{,}68$ passant-los primer tots a decimal i després tots a fracció.\par
\end{minipage}\par\vspace{1.6mm}
\par\noindent\begin{minipage}{\linewidth}
\textbf{5.} Situa aproximadament sobre una recta real els nombres $-2$, $1{,}5$, $\pi$ i $-\sqrt{2}$ (n'hi ha prou amb dues xifres decimals d'aproximació).\par
\end{minipage}\par\vspace{1.6mm}
\vspace{1mm}\hrule\vspace{3mm}

\section*{2. Intervals i valor absolut}

\par\noindent\begin{minipage}{\linewidth}
\textbf{6.} Escriu en forma d'interval, classifica'l (obert, tancat, semiobert, semirecta) i representa'l sobre la recta real:\par

\begin{itemize}[leftmargin=0pt,itemsep=3pt,topsep=3pt,parsep=0pt,label={}]
\item \textbf{a)} Nombres reals més grans que $2$ i menors o iguals que $7$.
\item \textbf{b)} Nombres reals menors que $-1$.
\item \textbf{c)} Nombres reals més grans o iguals que $0$.
\item \textbf{d)} Nombres reals entre $-3$ i $3$, ambdós inclosos.
\end{itemize}
\end{minipage}\par\vspace{1.6mm}
\par\noindent\begin{minipage}{\linewidth}
\textbf{7.} Per a cada interval, indica'n el tipus, escriu la desigualtat equivalent i representa'l sobre la recta real:\par

\begin{tabular}{@{}p{.46\linewidth}p{.46\linewidth}@{}}
\textbf{a)} $\left[-3,5\right)$ & \textbf{b)} $\left(2,+\infty\right)$ \\[2pt]
\textbf{c)} $\left(-\infty,4\right]$ & \textbf{d)} $\left(0,1\right)$ \\[2pt]
\end{tabular}\par
\end{minipage}\par\vspace{1.6mm}
\par\noindent\begin{minipage}{\linewidth}
\textbf{8.} Calcula:\par

\begin{tabular}{@{}p{.46\linewidth}p{.46\linewidth}@{}}
\textbf{a)} $\lvert 7\rvert$ & \textbf{b)} $\lvert -7\rvert$ \\[2pt]
\textbf{c)} $\lvert 0\rvert$ & \textbf{d)} $\lvert 3-8\rvert$ \\[2pt]
\textbf{e)} $\lvert 8-3\rvert$ & \textbf{f)} $-\lvert -4\rvert$ \\[2pt]
\end{tabular}\par
\end{minipage}\par\vspace{1.6mm}
\par\noindent\begin{minipage}{\linewidth}
\textbf{9.} Calcula la distància entre cada parell de punts:\par

\begin{tabular}{@{}p{.46\linewidth}p{.46\linewidth}@{}}
\textbf{a)} $a=2,\ b=9$ & \textbf{b)} $a=-5,\ b=-1$ \\[2pt]
\textbf{c)} $a=4,\ b=-3$ & \textbf{d)} $a=6,\ b=-2$ \\[2pt]
\end{tabular}\par
\end{minipage}\par\vspace{1.6mm}
\par\noindent\begin{minipage}{\linewidth}
\textbf{10.} Troba l'interval $\left(e_1,e_2\right)$ format pels punts que estan a una distància menor que $r$ del punt $a$:\par

\begin{itemize}[leftmargin=0pt,itemsep=3pt,topsep=3pt,parsep=0pt,label={}]
\item \textbf{a)} $a=5,\ r=3$
\item \textbf{b)} $a=-2,\ r=4$
\item \textbf{c)} $a=0,\ r=1{,}5$
\end{itemize}
\end{minipage}\par\vspace{1.6mm}
\par\noindent\begin{minipage}{\linewidth}
\textbf{11.} Resol i escriu la solució com a interval o, quan calgui, com a unió d'intervals. Recorda que $d\left(x,a\right)=|x-a|$, de manera que les dues notacions són equivalents:\par

\begin{tabular}{@{}p{.46\linewidth}p{.46\linewidth}@{}}
\textbf{a)} $\lvert x-3\rvert<5$ & \textbf{b)} $d\left(x,-2\right)<4$ \\[2pt]
\textbf{c)} $\lvert x-1\rvert\leq 6$ & \textbf{d)} $d\left(x,-5\right)<2$ \\[2pt]
\textbf{e)} $\lvert x-2\rvert>3$ & \textbf{f)} $d\left(x,-1\right)\geq 4$ \\[2pt]
\textbf{g)} $d\left(x,4\right)>1$ & \textbf{h)} $\lvert x+3\rvert\geq 2$ \\[2pt]
\end{tabular}\par

\emph{(Als apartats e)–h) la solució és el complementari d'un entorn: els punts que queden \textbf{lluny} del centre, no a prop.)}\par
\end{minipage}\par\vspace{1.6mm}
\par\noindent\begin{minipage}{\linewidth}
\textbf{12.} Troba $a$ i $r$ tals que l'interval solució de $|x-a|<r$ (o $\leq$, segons el cas) sigui:\par

\begin{itemize}[leftmargin=0pt,itemsep=3pt,topsep=3pt,parsep=0pt,label={}]
\item \textbf{a)} $\left(1,9\right)$
\item \textbf{b)} $\left(-4,4\right)$
\item \textbf{c)} $\left[-3,11\right]$
\end{itemize}
\end{minipage}\par\vspace{1.6mm}
\par\noindent\begin{minipage}{\linewidth}
\textbf{13.} Resol $|2x-6|<10$. \emph{(Indicació: pots treure factor comú $2$ dins el valor absolut abans d'aïllar.)}\par
\end{minipage}\par\vspace{1.6mm}
\par\noindent\begin{minipage}{\linewidth}
\textbf{14.} Resol $|3x+9|\leq 21$.\par
\end{minipage}\par\vspace{1.6mm}
\par\noindent\begin{minipage}{\linewidth}
\textbf{15.} Resol les inequacions lineals següents (sense valor absolut) i expressa la solució com a interval:\par

\begin{tabular}{@{}p{.46\linewidth}p{.46\linewidth}@{}}
\textbf{a)} $3x-4>11$ & \textbf{b)} $-2x+7\leq 1$ \\[2pt]
\textbf{c)} $5x-3<7$ & \textbf{d)} $-4x-1\geq 11$ \\[2pt]
\end{tabular}\par
\end{minipage}\par\vspace{1.6mm}
\par\noindent\begin{minipage}{\linewidth}
\textbf{16.} Resol:\par

\begin{itemize}[leftmargin=0pt,itemsep=3pt,topsep=3pt,parsep=0pt,label={}]
\item \textbf{a)} $\dfrac{x}{2}-3>1$
\item \textbf{b)} $2\left(x-1\right)\leq 3x+4$
\item \textbf{c)} $\dfrac{2x-1}{3}\geq 5$
\end{itemize}
\end{minipage}\par\vspace{1.6mm}
\par\noindent\begin{minipage}{\linewidth}
\textbf{17.} \emph{(aplicat)} La temperatura $T$ (en °C) d'una sala es manté dins un marge de $\pm 1{,}5$ °C respecte al valor de consigna de $21$ °C. Escriu aquesta condició amb un valor absolut i troba l'interval de temperatures possibles.\par
\end{minipage}\par\vspace{1.6mm}
\par\noindent\begin{minipage}{\linewidth}
\textbf{18.} \emph{(aplicat)} Una peça mecànica ha de tenir una longitud de $50\,\text{mm}$ amb una tolerància màxima de $0{,}2\,\text{mm}$.\par

\begin{itemize}[leftmargin=0pt,itemsep=3pt,topsep=3pt,parsep=0pt,label={}]
\item \textbf{a)} Escriu la condició amb valor absolut que ha de complir la longitud $\ell$.
\item \textbf{b)} Troba l'interval de longituds acceptades.
\end{itemize}
\end{minipage}\par\vspace{1.6mm}
\par\noindent\begin{minipage}{\linewidth}
\textbf{19.} \emph{(aplicat)} Dos vaixells naveguen per una mateixa ruta rectilínia; les seves posicions respecte a un far són $p_1=12\,\text{km}$ i $p_2=-5\,\text{km}$ (negatiu si són a l'altre costat del far).\par

\begin{itemize}[leftmargin=0pt,itemsep=3pt,topsep=3pt,parsep=0pt,label={}]
\item \textbf{a)} Calcula la distància entre els dos vaixells.
\item \textbf{b)} Una boia està ancorada a la posició $p=3\,\text{km}$. Un tercer vaixell s'ha de mantenir a una distància màxima de $8\,\text{km}$ de la boia: escriu la condició amb valor absolut i troba l'interval de posicions on es pot trobar.
\end{itemize}
\end{minipage}\par\vspace{1.6mm}
\par\noindent\begin{minipage}{\linewidth}
\textbf{20.} \emph{(aplicat)} En una carretera recta, el punt quilomètric $0$ marca l'entrada d'un poble i una antena de telefonia està situada al punt $p=7\,\text{km}$. La normativa obliga que qualsevol habitatge estigui a \textbf{més} de $2\,\text{km}$ de l'antena. (Considerem la carretera com una recta il·limitada.)\par

\begin{itemize}[leftmargin=0pt,itemsep=3pt,topsep=3pt,parsep=0pt,label={}]
\item \textbf{a)} Escriu amb un valor absolut la condició que ha de complir la posició $x$ d'un habitatge.
\item \textbf{b)} Troba el conjunt de posicions permeses i expressa'l com a unió d'intervals.
\item \textbf{c)} Un habitatge construït al punt quilomètric $8{,}5$ compleix la normativa?
\end{itemize}
\end{minipage}\par\vspace{1.6mm}
\par\noindent\begin{minipage}{\linewidth}
\textbf{21.} Troba tots els $x\in\mathbb{R}$ que compleixen alhora $|x-4|<6$ i $x>1$. Expressa la solució com a interval (intersecció).\par
\end{minipage}\par\vspace{1.6mm}
\par\noindent\begin{minipage}{\linewidth}
\textbf{22.} Escriu com un sol interval el conjunt de nombres reals que compleixen $-3\leq x<5$ o bé $4<x\leq 9$ (unió). Simplifica el resultat.\par
\end{minipage}\par\vspace{1.6mm}
\par\noindent\begin{minipage}{\linewidth}
\textbf{23.} Donats els conjunts $A$ i $B$, calcula $A\cap B$ i $A\cup B$, i expressa el resultat com a interval (o unió d'intervals) sempre que es pugui:\par

\begin{tabular}{@{}p{.46\linewidth}p{.46\linewidth}@{}}
\textbf{a)} $A=\left[-2,5\right)$, $B=\left(1,8\right]$ & \textbf{b)} $A=\left(-\infty,3\right]$, $B=\left(0,+\infty\right)$ \\[2pt]
\textbf{c)} $A=\left[-4,-1\right]$, $B=\left[2,6\right]$ & \textbf{d)} $A=\left(-3,4\right)$, $B=\left[0,+\infty\right)$ \\[2pt]
\end{tabular}\par
\end{minipage}\par\vspace{1.6mm}
\par\noindent\begin{minipage}{\linewidth}
\textbf{24.} Troba el conjunt de nombres reals que compleixen alhora les dues condicions\par

$$|x-1|<4 \qquad\text{i}\qquad |x+2|>1$$\par

Resol cada inequació per separat i després interseca les dues solucions.\par
\end{minipage}\par\vspace{1.6mm}
\vspace{1mm}\hrule\vspace{3mm}

\section*{3. Aproximacions i errors}

\par\noindent\begin{minipage}{\linewidth}
\textbf{25.} Sigui $e=2{,}718281828\ldots$ el nombre d'Euler. Trunca'l i arrodoneix-lo a $2$ i a $4$ xifres decimals.\par
\end{minipage}\par\vspace{1.6mm}
\par\noindent\begin{minipage}{\linewidth}
\textbf{26.} Indica el nombre de xifres significatives de:\par

\begin{itemize}[leftmargin=0pt,itemsep=3pt,topsep=3pt,parsep=0pt,label={}]
\item \textbf{a)} $0{,}0072$
\item \textbf{b)} $308{,}40$
\item \textbf{c)} $1{,}0500$
\item \textbf{d)} $6000$ (suposa que ve de mesurar amb precisió de milers)
\end{itemize}
\end{minipage}\par\vspace{1.6mm}
\par\noindent\begin{minipage}{\linewidth}
\textbf{27.} El valor exacte d'una magnitud és $x=12{,}763$ i s'aproxima per $x'=12{,}8$. Calcula l'error absolut i l'error relatiu (en tant per cent, arrodonit a $2$ xifres decimals).\par
\end{minipage}\par\vspace{1.6mm}
\par\noindent\begin{minipage}{\linewidth}
\textbf{28.} Un topògraf mesura la longitud d'un pont i obté $x'=248\,\text{m}$, mentre que el valor exacte és $x=250\,\text{m}$.\par

\begin{itemize}[leftmargin=0pt,itemsep=3pt,topsep=3pt,parsep=0pt,label={}]
\item \textbf{a)} Calcula l'error absolut.
\item \textbf{b)} Calcula l'error relatiu en tant per cent.
\end{itemize}
\end{minipage}\par\vspace{1.6mm}
\par\noindent\begin{minipage}{\linewidth}
\textbf{29.} \emph{(aplicat)} Es mesuren dues longituds amb el mateix error absolut d'$1\,\text{cm}$: la llargada d'un llapis ($15\,\text{cm}$) i la d'un camp de futbol ($100\,\text{m}$). Calcula l'error relatiu en cada cas i comenta quina mesura és relativament més precisa.\par
\end{minipage}\par\vspace{1.6mm}
\par\noindent\begin{minipage}{\linewidth}
\textbf{30.} Un nombre s'ha arrodonit a les centèsimes i s'ha obtingut $x'=7{,}42$. Entre quins valors es pot trobar el valor exacte $x$? Expressa la resposta en forma d'interval i vigila quins extrems s'inclouen i quins no.\par
\end{minipage}\par\vspace{1.6mm}
\par\noindent\begin{minipage}{\linewidth}
\textbf{31.} Cada valor següent s'ha obtingut per arrodoniment. Dona una fita de l'error absolut i una fita de l'error relatiu (en tant per cent, amb dues xifres decimals):\par

\begin{tabular}{@{}p{.46\linewidth}p{.46\linewidth}@{}}
\textbf{a)} $x'=3{,}25$ (a les centèsimes) & \textbf{b)} $x'=48{,}6$ (a les dècimes) \\[2pt]
\textbf{c)} $x'=1200$ (a les desenes) & \textbf{d)} $x'=0{,}038$ (a les mil·lèsimes) \\[2pt]
\end{tabular}\par
\end{minipage}\par\vspace{1.6mm}
\par\noindent\begin{minipage}{\linewidth}
\textbf{32.} Volem aproximar $\pi=3{,}14159265\ldots$ per arrodoniment, de manera que l'error relatiu sigui inferior al $0{,}01\%$. Quantes xifres decimals cal conservar com a mínim? Justifica-ho amb la fita de l'error i comprova després què passa si en conserves una menys.\par
\end{minipage}\par\vspace{1.6mm}
\par\noindent\begin{minipage}{\linewidth}
\textbf{33.} \emph{(aplicat)} La velocitat de la llum és $c=299\,792\,458\,\text{m/s}$, que sovint s'aproxima per $c'=300\,000\,000\,\text{m/s}$. Calcula l'error relatiu comès, en tant per cent (arrodoneix a $2$ xifres decimals). \emph{(Enllaça amb la notació científica que ve tot seguit.)}\par
\end{minipage}\par\vspace{1.6mm}
\vspace{1mm}\hrule\vspace{3mm}

\section*{4. Potències}

\par\noindent\begin{minipage}{\linewidth}
\textbf{34.} Calcula:\par

\begin{tabular}{@{}p{.46\linewidth}p{.46\linewidth}@{}}
\textbf{a)} $2^5$ & \textbf{b)} $\left(-3\right)^4$ \\[2pt]
\textbf{c)} $-3^4$ & \textbf{d)} $5^0$ \\[2pt]
\textbf{e)} $10^{-2}$ & \textbf{f)} $\left(\dfrac{2}{3}\right)^2$ \\[2pt]
\end{tabular}\par
\end{minipage}\par\vspace{1.6mm}
\par\noindent\begin{minipage}{\linewidth}
\textbf{35.} Aplica les propietats de les potències per simplificar (deixa el resultat com una única potència):\par

\begin{tabular}{@{}p{.46\linewidth}p{.46\linewidth}@{}}
\textbf{a)} $3^4\cdot 3^2$ & \textbf{b)} $\dfrac{5^7}{5^3}$ \\[2pt]
\textbf{c)} $\left(2^3\right)^2$ & \textbf{d)} $7^2\cdot 7^{-5}$ \\[2pt]
\textbf{e)} $\dfrac{4^{-1}}{4^{2}}$ &  \\[2pt]
\end{tabular}\par
\end{minipage}\par\vspace{1.6mm}
\par\noindent\begin{minipage}{\linewidth}
\textbf{36.} Calcula el valor numèric de cada apartat de l'exercici anterior.\par
\end{minipage}\par\vspace{1.6mm}
\par\noindent\begin{minipage}{\linewidth}
\textbf{37.} Simplifica i calcula el resultat:\par

\begin{tabular}{@{}p{.46\linewidth}p{.46\linewidth}@{}}
\textbf{a)} $\left(2\cdot 3\right)^2$ & \textbf{b)} $\left(\dfrac{3}{5}\right)^{-2}$ \\[2pt]
\textbf{c)} $2^3\cdot 2^{-1}\cdot 2^2$ & \textbf{d)} $\dfrac{\left(-2\right)^5}{\left(-2\right)^2}$ \\[2pt]
\end{tabular}\par
\end{minipage}\par\vspace{1.6mm}
\par\noindent\begin{minipage}{\linewidth}
\textbf{38.} Simplifica $\dfrac{2^5\cdot 4^2}{8^3}$ escrivint-ho tot amb la mateixa base, i calcula'n el valor.\par
\end{minipage}\par\vspace{1.6mm}
\par\noindent\begin{minipage}{\linewidth}
\textbf{39.} Simplifica $\dfrac{3^4\cdot 9}{27^2}$ escrivint-ho tot amb base $3$.\par
\end{minipage}\par\vspace{1.6mm}
\par\noindent\begin{minipage}{\linewidth}
\textbf{40.} Simplifica, amb $x\neq 0$: $\left(\dfrac{2x}{3}\right)^3 \cdot \left(\dfrac{3}{2x}\right)^2$.\par
\end{minipage}\par\vspace{1.6mm}
\par\noindent\begin{minipage}{\linewidth}
\textbf{41.} Simplifica combinant les propietats: $\dfrac{\left(2^3\right)^2\cdot 2^{-4}}{2^5\cdot 2^{-7}}$.\par
\end{minipage}\par\vspace{1.6mm}
\par\noindent\begin{minipage}{\linewidth}
\textbf{42.} Simplifica cada expressió com una única potència (totes les variables són no nul·les):\par

\begin{tabular}{@{}p{.46\linewidth}p{.46\linewidth}@{}}
\textbf{a)} $\dfrac{a^5\cdot a^{-2}}{a^{-3}\cdot a^4}$ & \textbf{b)} $\dfrac{\left(x^2\right)^3\cdot x^{-1}}{x^4}$ \\[2pt]
\textbf{c)} $\left(\dfrac{x^3}{y^2}\right)^2\cdot\dfrac{y^3}{x^4}$ & \textbf{d)} $\dfrac{\left(2a^2\right)^3}{4a^5}$ \\[2pt]
\end{tabular}\par
\end{minipage}\par\vspace{1.6mm}
\par\noindent\begin{minipage}{\linewidth}
\textbf{43.} Escriu en notació científica:\par

\begin{tabular}{@{}p{.46\linewidth}p{.46\linewidth}@{}}
\textbf{a)} $53\,000$ & \textbf{b)} $0{,}00048$ \\[2pt]
\textbf{c)} $7\,200\,000$ & \textbf{d)} $0{,}0000006$ \\[2pt]
\end{tabular}\par
\end{minipage}\par\vspace{1.6mm}
\par\noindent\begin{minipage}{\linewidth}
\textbf{44.} Escriu en forma decimal ordinària:\par

\begin{tabular}{@{}p{.46\linewidth}p{.46\linewidth}@{}}
\textbf{a)} $3{,}2\cdot 10^5$ & \textbf{b)} $9\cdot 10^{-3}$ \\[2pt]
\textbf{c)} $1{,}05\cdot 10^{2}$ & \textbf{d)} $6{,}7\cdot 10^{-6}$ \\[2pt]
\end{tabular}\par
\end{minipage}\par\vspace{1.6mm}
\par\noindent\begin{minipage}{\linewidth}
\textbf{45.} Calcula i expressa el resultat en notació científica:\par

\begin{tabular}{@{}p{.46\linewidth}p{.46\linewidth}@{}}
\textbf{a)} $\left(3\cdot10^4\right)\cdot\left(2\cdot10^3\right)$ & \textbf{b)} $\left(6\cdot10^{-2}\right)\cdot\left(5\cdot10^{5}\right)$ \\[2pt]
\textbf{c)} $\dfrac{8\cdot10^7}{2\cdot10^3}$ & \textbf{d)} $\dfrac{9\cdot10^{-2}}{3\cdot10^{4}}$ \\[2pt]
\end{tabular}\par
\end{minipage}\par\vspace{1.6mm}
\par\noindent\begin{minipage}{\linewidth}
\textbf{46.} Calcula, expressant el resultat en notació científica correcta: $\dfrac{\left(4\cdot10^{-3}\right)\cdot\left(6\cdot10^{8}\right)}{3\cdot10^{2}}$.\par
\end{minipage}\par\vspace{1.6mm}
\par\noindent\begin{minipage}{\linewidth}
\textbf{47.} \emph{(aplicat)} Un virus té un diàmetre aproximat de $1\cdot10^{-7}\,\text{m}$. Quants virus, un al costat de l'altre, calen per completar $1\,\text{cm}$?\par
\end{minipage}\par\vspace{1.6mm}
\par\noindent\begin{minipage}{\linewidth}
\textbf{48.} \emph{(aplicat)} La massa del Sol és aproximadament $1{,}8\cdot10^{30}\,\text{kg}$ i la de la Terra, $6\cdot10^{24}\,\text{kg}$. Quantes vegades més gran és la massa del Sol respecte a la de la Terra? Expressa el resultat en notació científica.\par
\end{minipage}\par\vspace{1.6mm}
\vspace{1mm}\hrule\vspace{3mm}

\section*{5. Arrels i radicals}

\par\noindent\begin{minipage}{\linewidth}
\textbf{49.} Calcula, si existeix, cada arrel. Si no existeix a $\mathbb{R}$, digues-ho:\par

\begin{tabular}{@{}p{.46\linewidth}p{.46\linewidth}@{}}
\textbf{a)} $\sqrt{25}$ & \textbf{b)} $\sqrt[3]{-27}$ \\[2pt]
\textbf{c)} $\sqrt{-9}$ & \textbf{d)} $\sqrt[4]{16}$ \\[2pt]
\textbf{e)} $\sqrt[5]{-32}$ & \textbf{f)} $\sqrt[3]{0}$ \\[2pt]
\end{tabular}\par
\end{minipage}\par\vspace{1.6mm}
\par\noindent\begin{minipage}{\linewidth}
\textbf{50.} Digues quantes arrels reals té cada equació i escriu-les:\par

\begin{tabular}{@{}p{.46\linewidth}p{.46\linewidth}@{}}
\textbf{a)} $x^2=49$ & \textbf{b)} $x^3=-8$ \\[2pt]
\textbf{c)} $x^2=-4$ & \textbf{d)} $x^4=81$ \\[2pt]
\end{tabular}\par
\end{minipage}\par\vspace{1.6mm}
\par\noindent\begin{minipage}{\linewidth}
\textbf{51.} Escriu cada radical com una potència d'exponent fraccionari:\par

\begin{tabular}{@{}p{.46\linewidth}p{.46\linewidth}@{}}
\textbf{a)} $\sqrt[3]{5}$ & \textbf{b)} $\sqrt[4]{2^3}$ \\[2pt]
\textbf{c)} $\sqrt{7}$ & \textbf{d)} $\sqrt[5]{x^2}$ \\[2pt]
\end{tabular}\par
\end{minipage}\par\vspace{1.6mm}
\par\noindent\begin{minipage}{\linewidth}
\textbf{52.} Escriu cada potència com un radical:\par

\begin{tabular}{@{}p{.46\linewidth}p{.46\linewidth}@{}}
\textbf{a)} $3^{1/2}$ & \textbf{b)} $2^{2/3}$ \\[2pt]
\textbf{c)} $x^{3/4}$ & \textbf{d)} $5^{1/4}$ \\[2pt]
\end{tabular}\par
\end{minipage}\par\vspace{1.6mm}
\par\noindent\begin{minipage}{\linewidth}
\textbf{53.} Simplifica reduint l'índex tant com sigui possible:\par

\begin{tabular}{@{}p{.46\linewidth}p{.46\linewidth}@{}}
\textbf{a)} $\sqrt[6]{4}$ & \textbf{b)} $\sqrt[4]{9}$ \\[2pt]
\textbf{c)} $\sqrt[8]{16}$ & \textbf{d)} $\sqrt[6]{27}$ \\[2pt]
\end{tabular}\par
\end{minipage}\par\vspace{1.6mm}
\par\noindent\begin{minipage}{\linewidth}
\textbf{54.} Extreu tots els factors possibles de cada radical (suposa $x>0$, $a>0$):\par

\begin{tabular}{@{}p{.46\linewidth}p{.46\linewidth}@{}}
\textbf{a)} $\sqrt{75}$ & \textbf{b)} $\sqrt[3]{54}$ \\[2pt]
\textbf{c)} $\sqrt{200}$ & \textbf{d)} $\sqrt[3]{128}$ \\[2pt]
\textbf{e)} $\sqrt{x^5}$ & \textbf{f)} $\sqrt[3]{8x^4}$ \\[2pt]
\textbf{g)} $\sqrt[4]{\dfrac{x^8}{16}}$ & \textbf{h)} $\sqrt{18a^3}$ \\[2pt]
\end{tabular}\par
\end{minipage}\par\vspace{1.6mm}
\par\noindent\begin{minipage}{\linewidth}
\textbf{55.} Introdueix els factors dins del radical i simplifica (suposa $x>0$, $a>0$):\par

\begin{tabular}{@{}p{.46\linewidth}p{.46\linewidth}@{}}
\textbf{a)} $3\sqrt{2}$ & \textbf{b)} $2\sqrt[3]{5}$ \\[2pt]
\textbf{c)} $x\sqrt{x}$ & \textbf{d)} $4\sqrt[3]{2}$ \\[2pt]
\textbf{e)} $\dfrac{2}{a}\sqrt{\dfrac{a^3}{4}}$ & \textbf{f)} $\dfrac{1}{2}\sqrt{20}$ \\[2pt]
\end{tabular}\par
\end{minipage}\par\vspace{1.6mm}
\par\noindent\begin{minipage}{\linewidth}
\textbf{56.} Opera primer el radicand i després extreu els factors que puguis:\par

\begin{tabular}{@{}p{.46\linewidth}p{.46\linewidth}@{}}
\textbf{a)} $\sqrt{\dfrac{1}{4}+\dfrac{1}{9}}$ & \textbf{b)} $\sqrt{\dfrac{x}{4}+2x}$, amb $x\geq 0$ \\[2pt]
\textbf{c)} $\sqrt[3]{27x^2+27}$ & \textbf{d)} $\sqrt{4x^2+4x+1}$ \\[2pt]
\end{tabular}\par
\end{minipage}\par\vspace{1.6mm}
\par\noindent\begin{minipage}{\linewidth}
\textbf{57.} Sense calculadora, ordena de més petit a més gran, passant-los prèviament a índex comú: $\sqrt{2}$, $\sqrt[3]{3}$, $\sqrt[6]{5}$, $\sqrt[3]{4}$.\par
\end{minipage}\par\vspace{1.6mm}
\par\noindent\begin{minipage}{\linewidth}
\textbf{58.} Multiplica i simplifica:\par

\begin{tabular}{@{}p{.46\linewidth}p{.46\linewidth}@{}}
\textbf{a)} $\sqrt{3}\cdot\sqrt{12}$ & \textbf{b)} $\sqrt[3]{2}\cdot\sqrt[3]{4}$ \\[2pt]
\textbf{c)} $\sqrt{5}\cdot\sqrt{20}$ & \textbf{d)} $\sqrt[4]{2}\cdot\sqrt[4]{8}$ \\[2pt]
\end{tabular}\par
\end{minipage}\par\vspace{1.6mm}
\par\noindent\begin{minipage}{\linewidth}
\textbf{59.} Divideix i simplifica:\par

\begin{tabular}{@{}p{.46\linewidth}p{.46\linewidth}@{}}
\textbf{a)} $\dfrac{\sqrt{50}}{\sqrt{2}}$ & \textbf{b)} $\dfrac{\sqrt[3]{54}}{\sqrt[3]{2}}$ \\[2pt]
\textbf{c)} $\dfrac{\sqrt{45}}{\sqrt{5}}$ & \textbf{d)} $\dfrac{\sqrt[4]{48}}{\sqrt[4]{3}}$ \\[2pt]
\end{tabular}\par
\end{minipage}\par\vspace{1.6mm}
\par\noindent\begin{minipage}{\linewidth}
\textbf{60.} Opera i simplifica, portant primer els radicals a índex comú:\par

\begin{itemize}[leftmargin=0pt,itemsep=3pt,topsep=3pt,parsep=0pt,label={}]
\item \textbf{a)} $\sqrt{2}\cdot\sqrt[3]{2}$
\item \textbf{b)} $\dfrac{\sqrt[3]{9}}{\sqrt{3}}$
\end{itemize}
\end{minipage}\par\vspace{1.6mm}
\par\noindent\begin{minipage}{\linewidth}
\textbf{61.} Simplifica:\par

\begin{tabular}{@{}p{.46\linewidth}p{.46\linewidth}@{}}
\textbf{a)} $\sqrt{\sqrt{16}}$ & \textbf{b)} $\sqrt[3]{\sqrt[4]{x}}$ \\[2pt]
\textbf{c)} $\left(\sqrt[3]{5}\right)^2$ & \textbf{d)} $\left(\sqrt{7}\right)^4$ \\[2pt]
\end{tabular}\par
\end{minipage}\par\vspace{1.6mm}
\par\noindent\begin{minipage}{\linewidth}
\textbf{62.} Expressa cada expressió com un sol radical i simplifica:\par

\begin{tabular}{@{}p{.46\linewidth}p{.46\linewidth}@{}}
\textbf{a)} $\sqrt{5\sqrt{5}}$ & \textbf{b)} $\sqrt[3]{2\sqrt{2}}$ \\[2pt]
\textbf{c)} $\dfrac{\left(\sqrt[3]{16}\right)^2}{\sqrt{32}}$ & \textbf{d)} $\sqrt{3\sqrt[3]{3}}$ \\[2pt]
\end{tabular}\par
\end{minipage}\par\vspace{1.6mm}
\par\noindent\begin{minipage}{\linewidth}
\textbf{63.} Suma o resta els radicals (ja són semblants):\par

\begin{itemize}[leftmargin=0pt,itemsep=3pt,topsep=3pt,parsep=0pt,label={}]
\item \textbf{a)} $3\sqrt2+5\sqrt2$
\item \textbf{b)} $7\sqrt[3]{3}-2\sqrt[3]{3}$
\item \textbf{c)} $2\sqrt5-6\sqrt5+\sqrt5$
\end{itemize}
\end{minipage}\par\vspace{1.6mm}
\par\noindent\begin{minipage}{\linewidth}
\textbf{64.} Extreu factors per identificar els radicals semblants i suma'ls:\par

\begin{itemize}[leftmargin=0pt,itemsep=3pt,topsep=3pt,parsep=0pt,label={}]
\item \textbf{a)} $\sqrt{8}+\sqrt{18}$
\item \textbf{b)} $\sqrt{50}-\sqrt{32}$
\item \textbf{c)} $2\sqrt{12}+3\sqrt{27}$
\end{itemize}
\end{minipage}\par\vspace{1.6mm}
\par\noindent\begin{minipage}{\linewidth}
\textbf{65.} Calcula i simplifica (suposa $a>0$, $x>0$):\par

\begin{itemize}[leftmargin=0pt,itemsep=3pt,topsep=3pt,parsep=0pt,label={}]
\item \textbf{a)} $3\sqrt{20}-2\sqrt{45}+\sqrt{80}$
\item \textbf{b)} $4\sqrt[3]{8a}-\sqrt[3]{27a}$
\item \textbf{c)} $\sqrt{9x}+\sqrt{25x}-\sqrt{4x}$
\end{itemize}
\end{minipage}\par\vspace{1.6mm}
\par\noindent\begin{minipage}{\linewidth}
\textbf{66.} Racionalitza el denominador:\par

\begin{tabular}{@{}p{.46\linewidth}p{.46\linewidth}@{}}
\textbf{a)} $\dfrac{5}{\sqrt3}$ & \textbf{b)} $\dfrac{2}{\sqrt5}$ \\[2pt]
\textbf{c)} $\dfrac{6}{\sqrt2}$ & \textbf{d)} $\dfrac{1}{\sqrt7}$ \\[2pt]
\end{tabular}\par
\end{minipage}\par\vspace{1.6mm}
\par\noindent\begin{minipage}{\linewidth}
\textbf{67.} Racionalitza:\par

\begin{itemize}[leftmargin=0pt,itemsep=3pt,topsep=3pt,parsep=0pt,label={}]
\item \textbf{a)} $\dfrac{4}{2\sqrt3}$
\item \textbf{b)} $\dfrac{3}{\sqrt{12}}$
\item \textbf{c)} $\dfrac{10}{3\sqrt5}$
\end{itemize}
\end{minipage}\par\vspace{1.6mm}
\par\noindent\begin{minipage}{\linewidth}
\textbf{68.} Racionalitza utilitzant el conjugat:\par

\begin{itemize}[leftmargin=0pt,itemsep=3pt,topsep=3pt,parsep=0pt,label={}]
\item \textbf{a)} $\dfrac{1}{\sqrt5-2}$
\item \textbf{b)} $\dfrac{3}{\sqrt7+1}$
\item \textbf{c)} $\dfrac{2}{3-\sqrt2}$
\end{itemize}
\end{minipage}\par\vspace{1.6mm}
\par\noindent\begin{minipage}{\linewidth}
\textbf{69.} Racionalitza i simplifica:\par

\begin{itemize}[leftmargin=0pt,itemsep=3pt,topsep=3pt,parsep=0pt,label={}]
\item \textbf{a)} $\dfrac{5}{\sqrt3-\sqrt2}$
\item \textbf{b)} $\dfrac{1}{\sqrt{n}}$, amb $n\in\mathbb{N}$
\item \textbf{c)} $\dfrac{a-b}{\sqrt{a}+\sqrt{b}}$, amb $a,b\in\mathbb{N}$ i $a\neq b$
\end{itemize}
\end{minipage}\par\vspace{1.6mm}
\par\noindent\begin{minipage}{\linewidth}
\textbf{70.} Simplifica combinant diverses propietats: $\sqrt{18}+\dfrac{4}{\sqrt2}-\sqrt{50}$.\par
\end{minipage}\par\vspace{1.6mm}
\par\noindent\begin{minipage}{\linewidth}
\textbf{71.} \emph{(aplicat)} Els catets d'un triangle rectangle fan $a=\sqrt{18}$ i $b=\sqrt{32}$. Calcula, simplificant el resultat, la hipotenusa $c=\sqrt{a^2+b^2}$.\par
\end{minipage}\par\vspace{1.6mm}
\par\noindent\begin{minipage}{\linewidth}
\textbf{72.} \emph{(aplicat)} El costat d'un quadrat mesura $\ell=3\sqrt2\,\text{cm}$. Calcula, simplificant:\par

\begin{itemize}[leftmargin=0pt,itemsep=3pt,topsep=3pt,parsep=0pt,label={}]
\item \textbf{a)} L'àrea del quadrat.
\item \textbf{b)} La diagonal (recorda que la diagonal d'un quadrat de costat $\ell$ és $\ell\sqrt2$).
\end{itemize}
\end{minipage}\par\vspace{1.6mm}
\par\noindent\begin{minipage}{\linewidth}
\textbf{73.} Comprova, amb $a=9$ i $b=4$, que $\sqrt{a+b}\neq\sqrt{a}+\sqrt{b}$ en general. Calcula els dos costats i compara.\par
\end{minipage}\par\vspace{1.6mm}
\par\noindent\begin{minipage}{\linewidth}
\textbf{74.} Simplifica $\dfrac{\sqrt{x^5}}{\sqrt{x}}$, amb $x>0$, sense radicals al denominador i amb l'exponent més senzill possible.\par
\end{minipage}\par\vspace{1.6mm}
\vspace{1mm}\hrule\vspace{3mm}

\section*{6. Logaritmes}

\subsection*{6.1 Definició i càlcul directe}

\par\noindent\begin{minipage}{\linewidth}
\textbf{75.} Escriu en forma de logaritme:\par

\begin{tabular}{@{}p{.46\linewidth}p{.46\linewidth}@{}}
\textbf{a)} $2^5=32$ & \textbf{b)} $10^3=1000$ \\[2pt]
\textbf{c)} $5^0=1$ & \textbf{d)} $3^{-2}=\dfrac{1}{9}$ \\[2pt]
\end{tabular}\par
\end{minipage}\par\vspace{1.6mm}
\par\noindent\begin{minipage}{\linewidth}
\textbf{76.} Escriu en forma de potència:\par

\begin{tabular}{@{}p{.46\linewidth}p{.46\linewidth}@{}}
\textbf{a)} $\log_2 16=4$ & \textbf{b)} $\log 100=2$ \\[2pt]
\textbf{c)} $\log_3 1=0$ & \textbf{d)} $\log_5 5=1$ \\[2pt]
\end{tabular}\par
\end{minipage}\par\vspace{1.6mm}
\par\noindent\begin{minipage}{\linewidth}
\textbf{77.} Calcula sense calculadora:\par

\begin{tabular}{@{}p{.46\linewidth}p{.46\linewidth}@{}}
\textbf{a)} $\log_2 8$ & \textbf{b)} $\log_3 81$ \\[2pt]
\textbf{c)} $\log 10000$ & \textbf{d)} $\log_5 1$ \\[2pt]
\textbf{e)} $\log_7 7$ & \textbf{f)} $\ln e$ \\[2pt]
\end{tabular}\par
\end{minipage}\par\vspace{1.6mm}
\par\noindent\begin{minipage}{\linewidth}
\textbf{78.} Calcula sense calculadora, escrivint prèviament el nombre com una potència de la base:\par

\begin{tabular}{@{}p{.46\linewidth}p{.46\linewidth}@{}}
\textbf{a)} $\log_2 \dfrac{1}{4}$ & \textbf{b)} $\log_3 \dfrac{1}{27}$ \\[2pt]
\textbf{c)} $\log_5 \sqrt5$ & \textbf{d)} $\log_2 \sqrt[3]{4}$ \\[2pt]
\end{tabular}\par
\end{minipage}\par\vspace{1.6mm}
\subsection*{6.2 Propietats amb valors numèrics}

\par\noindent\begin{minipage}{\linewidth}
\textbf{79.} Aplica la propietat del logaritme d'un producte per desenvolupar:\par

\begin{itemize}[leftmargin=0pt,itemsep=3pt,topsep=3pt,parsep=0pt,label={}]
\item \textbf{a)} $\log_2\left(4\cdot 8\right)$
\item \textbf{b)} $\log\left(5\cdot 20\right)$
\item \textbf{c)} $\log_3\left(x\cdot y\right)$, amb $x,y>0$
\end{itemize}
\end{minipage}\par\vspace{1.6mm}
\par\noindent\begin{minipage}{\linewidth}
\textbf{80.} Aplica la propietat del logaritme d'un quocient:\par

\begin{itemize}[leftmargin=0pt,itemsep=3pt,topsep=3pt,parsep=0pt,label={}]
\item \textbf{a)} $\log_2 \dfrac{32}{4}$
\item \textbf{b)} $\log \dfrac{1000}{10}$
\item \textbf{c)} $\log_5 \dfrac{x}{25}$, amb $x>0$
\end{itemize}
\end{minipage}\par\vspace{1.6mm}
\par\noindent\begin{minipage}{\linewidth}
\textbf{81.} Aplica la propietat de la potència:\par

\begin{tabular}{@{}p{.46\linewidth}p{.46\linewidth}@{}}
\textbf{a)} $\log_2 8^5$ & \textbf{b)} $\log_3 9^{-2}$ \\[2pt]
\textbf{c)} $\log\left(10^7\right)$ & \textbf{d)} $\log_a a^n$ \\[2pt]
\end{tabular}\par
\end{minipage}\par\vspace{1.6mm}
\par\noindent\begin{minipage}{\linewidth}
\textbf{82.} Aplica la propietat de l'arrel:\par

\begin{itemize}[leftmargin=0pt,itemsep=3pt,topsep=3pt,parsep=0pt,label={}]
\item \textbf{a)} $\log_2 \sqrt[3]{8}$
\item \textbf{b)} $\log_5\sqrt{25}$
\item \textbf{c)} $\log_3\sqrt[4]{81}$
\end{itemize}
\end{minipage}\par\vspace{1.6mm}
\par\noindent\begin{minipage}{\linewidth}
\textbf{83.} Desenvolupa aplicant totes les propietats necessàries: $\log_2\left(\dfrac{8\sqrt2}{4}\right)$.\par
\end{minipage}\par\vspace{1.6mm}
\par\noindent\begin{minipage}{\linewidth}
\textbf{84.} Desenvolupa: $\log\left(\dfrac{100\cdot\sqrt{10}}{1000}\right)$.\par
\end{minipage}\par\vspace{1.6mm}
\par\noindent\begin{minipage}{\linewidth}
\textbf{85.} Redueix a un sol logaritme:\par

\begin{tabular}{@{}p{.46\linewidth}p{.46\linewidth}@{}}
\textbf{a)} $\log_2 5+\log_2 3$ & \textbf{b)} $\log 20-\log 4$ \\[2pt]
\textbf{c)} $3\log_5 2$ & \textbf{d)} $\dfrac12 \log_3 9$ \\[2pt]
\end{tabular}\par
\end{minipage}\par\vspace{1.6mm}
\par\noindent\begin{minipage}{\linewidth}
\textbf{86.} Redueix a un sol logaritme i calcula'n el valor: $2\log_2 3-\log_2 \dfrac{9}{4}$.\par
\end{minipage}\par\vspace{1.6mm}
\par\noindent\begin{minipage}{\linewidth}
\textbf{87.} Utilitza el canvi de base per expressar en base $10$ (deixa'l indicat, sense calcular-lo):\par

\begin{itemize}[leftmargin=0pt,itemsep=3pt,topsep=3pt,parsep=0pt,label={}]
\item \textbf{a)} $\log_2 7$
\item \textbf{b)} $\log_5 12$
\end{itemize}
\end{minipage}\par\vspace{1.6mm}
\subsection*{6.3 Propietats amb expressions literals}

\begin{quote}\small En tots els exercicis d'aquest apartat, suposem $a>0$, $a\neq1$, i totes les variables ($x$, $y$, $z$) estrictament positives.\end{quote}

\par\noindent\begin{minipage}{\linewidth}
\textbf{88.} Desenvolupa aplicant les propietats dels logaritmes:\par

\begin{tabular}{@{}p{.46\linewidth}p{.46\linewidth}@{}}
\textbf{a)} $\log_a\left(x^3 y^2\right)$ & \textbf{b)} $\log_a \dfrac{x^2}{y^3}$ \\[2pt]
\textbf{c)} $\log_a \dfrac{xy}{z}$ & \textbf{d)} $\log_a \sqrt{x}$ \\[2pt]
\end{tabular}\par
\end{minipage}\par\vspace{1.6mm}
\par\noindent\begin{minipage}{\linewidth}
\textbf{89.} Desenvolupa completament, fins que no quedi cap producte, quocient, potència ni arrel dins d'un logaritme:\par

\begin{tabular}{@{}p{.46\linewidth}p{.46\linewidth}@{}}
\textbf{a)} $\log_a \dfrac{x^3\sqrt{y}}{z^2}$ & \textbf{b)} $\log_a \sqrt{\dfrac{x}{y}}$ \\[2pt]
\textbf{c)} $\log_a \dfrac{x^2\sqrt[3]{y}}{z}$ & \textbf{d)} $\log_a \left(x\sqrt{yz}\right)$ \\[2pt]
\end{tabular}\par
\end{minipage}\par\vspace{1.6mm}
\par\noindent\begin{minipage}{\linewidth}
\textbf{90.} Redueix a un sol logaritme:\par

\begin{tabular}{@{}p{.46\linewidth}p{.46\linewidth}@{}}
\textbf{a)} $2\log_a x+3\log_a y$ & \textbf{b)} $\log_a x-\dfrac12\log_a y$ \\[2pt]
\textbf{c)} $\log_a x+\log_a y-\log_a z$ & \textbf{d)} $4\log_a x$ \\[2pt]
\end{tabular}\par
\end{minipage}\par\vspace{1.6mm}
\par\noindent\begin{minipage}{\linewidth}
\textbf{91.} Redueix a un sol logaritme:\par

\begin{tabular}{@{}p{.46\linewidth}p{.46\linewidth}@{}}
\textbf{a)} $3\log_a x-\log_a y-2\log_a z$ & \textbf{b)} $\dfrac12\left(\log_a x+\log_a y\right)$ \\[2pt]
\textbf{c)} $2\log_a x-\dfrac13\log_a y$ & \textbf{d)} $\dfrac13\log_a x+\dfrac23\log_a y-\log_a z$ \\[2pt]
\end{tabular}\par
\end{minipage}\par\vspace{1.6mm}
\par\noindent\begin{minipage}{\linewidth}
\textbf{92.} Sabent que $\log_a 2=m$ i $\log_a 3=n$, expressa en funció de $m$ i $n$:\par

\begin{tabular}{@{}p{.46\linewidth}p{.46\linewidth}@{}}
\textbf{a)} $\log_a 6$ & \textbf{b)} $\log_a 12$ \\[2pt]
\textbf{c)} $\log_a \dfrac{3}{2}$ & \textbf{d)} $\log_a \sqrt{2}$ \\[2pt]
\textbf{e)} $\log_a 18$ & \textbf{f)} $\log_a \dfrac{8}{9}$ \\[2pt]
\end{tabular}\par
\end{minipage}\par\vspace{1.6mm}
\par\noindent\begin{minipage}{\linewidth}
\textbf{93.} Sabent que $\log 2 \approx 0{,}301$, calcula sense calculadora (recorda que $\log 10=1$):\par

\begin{tabular}{@{}p{.46\linewidth}p{.46\linewidth}@{}}
\textbf{a)} $\log 5$ & \textbf{b)} $\log 4$ \\[2pt]
\textbf{c)} $\log 50$ & \textbf{d)} $\log 0{,}2$ \\[2pt]
\textbf{e)} $\log \sqrt{2}$ & \textbf{f)} $\log 40$ \\[2pt]
\end{tabular}\par
\end{minipage}\par\vspace{1.6mm}
\par\noindent\begin{minipage}{\linewidth}
\textbf{94.} Anomenem $L=\log_a x$. Expressa en funció de $L$ i simplifica al màxim:\par

$$\log_a x^2-\log_a \sqrt{x}+\log_a \dfrac{1}{x}$$\par
\end{minipage}\par\vspace{1.6mm}
\par\noindent\begin{minipage}{\linewidth}
\textbf{95.} Demostra, utilitzant el canvi de base, que per a qualsevol $a,b>0$ amb $a\neq1$ i $b\neq1$ es compleix\par

$$\log_a b \cdot \log_b a = 1$$\par
\end{minipage}\par\vspace{1.6mm}
\par\noindent\begin{minipage}{\linewidth}
\textbf{96.} Demostra que $\log_{a^2} x = \dfrac12 \log_a x$. \emph{(Indicació: aplica el canvi de base a la base $a$.)}\par
\end{minipage}\par\vspace{1.6mm}
\subsection*{6.4 Aplicacions}

\par\noindent\begin{minipage}{\linewidth}
\textbf{97.} \emph{(aplicat)} Un capital de $C_0=5\,000$ \euro{} s'inverteix a un interès compost anual del $4\%$. Quants anys calen, aproximadament, perquè el capital arribi a $C_f=7\,000$ \euro{}? \emph{(Usa $C_f=C_0\left(1+i\right)^t$ i el logaritme per aïllar $t$.)}\par
\end{minipage}\par\vspace{1.6mm}
\par\noindent\begin{minipage}{\linewidth}
\textbf{98.} \emph{(aplicat)} La magnitud $M$ d'un terratrèmol (escala de Richter) es relaciona amb l'energia alliberada $E$ mitjançant $M=\dfrac{2}{3}\log\left(\dfrac{E}{E_0}\right)$, amb $E_0$ constant. Si un terratrèmol allibera $1000$ vegades més energia que un altre ($E=1000\,E_1$), quant més gran és la seva magnitud $M$ respecte a $M_1$? \emph{(Calcula $M-M_1$.)}\par
\end{minipage}\par\vspace{1.6mm}
\vspace{1mm}\hrule\vspace{3mm}

\section*{7. Problemes de repàs i aplicacions}

\par\noindent\begin{minipage}{\linewidth}
\textbf{99.} \emph{(repàs)} Simplifica combinant potències i arrels: $\sqrt{2^6}\cdot 2^{-1}$.\par
\end{minipage}\par\vspace{1.6mm}
\par\noindent\begin{minipage}{\linewidth}
\textbf{100.} \emph{(repàs)} Calcula, en notació científica: $\left(2\cdot10^3\right)^2 \cdot \left(5\cdot10^{-2}\right)$.\par
\end{minipage}\par\vspace{1.6mm}
\par\noindent\begin{minipage}{\linewidth}
\textbf{101.} \emph{(repàs)} Simplifica $\sqrt{\dfrac{4\cdot10^{-2}}{9}}$ i expressa el resultat com una fracció.\par
\end{minipage}\par\vspace{1.6mm}
\par\noindent\begin{minipage}{\linewidth}
\textbf{102.} \emph{(repàs)} Troba l'interval solució de $|x-3|<\sqrt{16}$.\par
\end{minipage}\par\vspace{1.6mm}
\par\noindent\begin{minipage}{\linewidth}
\textbf{103.} \emph{(repàs)} Calcula l'error relatiu, en tant per cent, de l'aproximació $\sqrt2 \approx 1{,}414$ (arrodoneix a $3$ xifres decimals).\par
\end{minipage}\par\vspace{1.6mm}
\par\noindent\begin{minipage}{\linewidth}
\textbf{104.} \emph{(repàs · aplicat)} Una població de bacteris es multiplica per $2$ cada hora, segons $N\left(t\right)=N_0\cdot 2^t$, amb $N_0=500$.\par

\begin{itemize}[leftmargin=0pt,itemsep=3pt,topsep=3pt,parsep=0pt,label={}]
\item \textbf{a)} Quants bacteris hi ha al cap de $5$ hores?
\item \textbf{b)} Quantes hores calen perquè la població arribi a $64\,000$ bacteris?
\end{itemize}
\end{minipage}\par\vspace{1.6mm}
\par\noindent\begin{minipage}{\linewidth}
\textbf{105.} \emph{(repàs · aplicat)} Un dipòsit cúbic té un volum de $V=125\,\text{m}^3$.\par

\begin{itemize}[leftmargin=0pt,itemsep=3pt,topsep=3pt,parsep=0pt,label={}]
\item \textbf{a)} Calcula, sense calculadora, l'aresta del cub ($V=a^3 \Rightarrow a=\sqrt[3]{V}$).
\item \textbf{b)} Si es vol doblar el volum, quina hauria de ser la nova aresta $a'$? Deixa-la com a radical simplificat.
\end{itemize}
\end{minipage}\par\vspace{1.6mm}
\par\noindent\begin{minipage}{\linewidth}
\textbf{106.} \emph{(repàs · aplicat)} La intensitat sonora $\beta$ (en decibels) es calcula com $\beta=10\log\left(\dfrac{I}{I_0}\right)$, amb $I_0=10^{-12}\,\text{W/m}^2$ el llindar d'audició. Calcula la intensitat sonora, en decibels, d'un so amb $I=10^{-4}\,\text{W/m}^2$.\par
\end{minipage}\par\vspace{1.6mm}
\par\noindent\begin{minipage}{\linewidth}
\textbf{107.} \emph{(repàs global)} Simplifica pas a pas fins a obtenir un únic nombre:\par

$$\dfrac{\sqrt{50}-\sqrt{8}}{\sqrt2} + \log_2 16 - |{-3}|$$\par
\end{minipage}\par\vspace{1.6mm}
\par\noindent\begin{minipage}{\linewidth}
\textbf{108.} \emph{(repàs global)} Considera $A=\dfrac{3^2\cdot\sqrt4}{3}$ i $B=\log_2\left(\dfrac{32}{4}\right)$.\par

\begin{itemize}[leftmargin=0pt,itemsep=3pt,topsep=3pt,parsep=0pt,label={}]
\item \textbf{a)} Calcula el valor numèric de $A$ i de $B$.
\item \textbf{b)} Troba l'interval solució de $|x-A|<B$.
\end{itemize}
\end{minipage}\par\vspace{1.6mm}
\par\noindent\begin{minipage}{\linewidth}
\textbf{109.} Digues si cada afirmació és certa o falsa. Si és falsa, dona un contraexemple:\par

\begin{itemize}[leftmargin=0pt,itemsep=3pt,topsep=3pt,parsep=0pt,label={}]
\item \textbf{a)} $\sqrt{a^2}=a$ per a tot $a\in\mathbb{R}$.
\item \textbf{b)} $|a+b|=|a|+|b|$ per a tot $a,b\in\mathbb{R}$.
\item \textbf{c)} $\log\left(x+y\right)=\log x+\log y$ per a tot $x,y>0$.
\item \textbf{d)} $\sqrt[3]{a}$ existeix per a tot $a\in\mathbb{R}$.
\item \textbf{e)} Si $|x|<3$, aleshores $x<3$.
\item \textbf{f)} $\left(a+b\right)^2=a^2+b^2$ per a tot $a,b\in\mathbb{R}$.
\end{itemize}
\end{minipage}\par\vspace{1.6mm}
\par\noindent\begin{minipage}{\linewidth}
\textbf{110.} \emph{(aplicat)} El període d'oscil·lació $T$ d'un pèndol simple de longitud $L$ ve donat per\par

$$T=2\pi\sqrt{\dfrac{L}{g}},$$\par

on $g$ és l'acceleració de la gravetat (constant). Demostra que, per aconseguir que el període es dupliqui, cal multiplicar la longitud del pèndol per $4$.\par
\end{minipage}\par\vspace{1.6mm}
\vspace{1mm}\hrule\vspace{3mm}

\section*{Solucions}
\small\begin{multicols}{2}


\textbf{Bloc 1.}\par
\par\noindent\begin{minipage}{\linewidth}
\textbf{1.} a) $\mathbb{Z},\mathbb{Q}$ · b) $\mathbb{Q}$ · c) $4$: $\mathbb{N},\mathbb{Z},\mathbb{Q}$ · d) $\mathbb{R}\setminus\mathbb{Q}$ · e) $\mathbb{N},\mathbb{Z},\mathbb{Q}$ · f) $\mathbb{Q}$ · g) $-3$: $\mathbb{Z},\mathbb{Q}$ · h) $\mathbb{R}\setminus\mathbb{Q}$.\par
\end{minipage}\par\vspace{1.6mm}
\par\noindent\begin{minipage}{\linewidth}
\textbf{2.} a) V · b) F · c) V · d) F ($0{,}\overline6=\tfrac23$).\par
\end{minipage}\par\vspace{1.6mm}
\par\noindent\begin{minipage}{\linewidth}
\textbf{3.} $-3,\ -1{,}5,\ 0,\ \tfrac72,\ 4$.\par
\end{minipage}\par\vspace{1.6mm}
\par\noindent\begin{minipage}{\linewidth}
\textbf{4.} $0{,}68$ (perquè $\tfrac23=0{,}\overline6<0{,}68$).\par

\textbf{Bloc 2.}\par
\end{minipage}\par\vspace{1.6mm}
\par\noindent\begin{minipage}{\linewidth}
\textbf{6.} a) $\left(2,7\right]$, semiobert · b) $\left(-\infty,-1\right)$, semirecta · c) $\left[0,+\infty\right)$, semirecta · d) $\left[-3,3\right]$, tancat.\par
\end{minipage}\par\vspace{1.6mm}
\par\noindent\begin{minipage}{\linewidth}
\textbf{7.} a) $-3\le x<5$ · b) $x>2$ · c) $x\le4$ · d) $0<x<1$.\par
\end{minipage}\par\vspace{1.6mm}
\par\noindent\begin{minipage}{\linewidth}
\textbf{8.} a) $7$ · b) $7$ · c) $0$ · d) $5$ · e) $5$ · f) $-4$.\par
\end{minipage}\par\vspace{1.6mm}
\par\noindent\begin{minipage}{\linewidth}
\textbf{9.} a) $7$ · b) $4$ · c) $7$ · d) $8$.\par
\end{minipage}\par\vspace{1.6mm}
\par\noindent\begin{minipage}{\linewidth}
\textbf{10.} a) $\left(2,8\right)$ · b) $\left(-6,2\right)$ · c) $\left(-1{,}5,\,1{,}5\right)$.\par
\end{minipage}\par\vspace{1.6mm}
\par\noindent\begin{minipage}{\linewidth}
\textbf{11.} a) $\left(-2,8\right)$ · b) $\left(-6,2\right)$ · c) $\left[-5,7\right]$ · d) $\left(-7,-3\right)$ · e) $\left(-\infty,-1\right)\cup\left(5,+\infty\right)$ · f) $\left(-\infty,-5\right]\cup\left[3,+\infty\right)$ · g) $\left(-\infty,3\right)\cup\left(5,+\infty\right)$ · h) $\left(-\infty,-5\right]\cup\left[-1,+\infty\right)$.\par
\end{minipage}\par\vspace{1.6mm}
\par\noindent\begin{minipage}{\linewidth}
\textbf{12.} a) $a=5,\ r=4$ · b) $a=0,\ r=4$ · c) $a=4,\ r=7$.\par
\end{minipage}\par\vspace{1.6mm}
\par\noindent\begin{minipage}{\linewidth}
\textbf{13.} $\left(-2,8\right)$.\par
\end{minipage}\par\vspace{1.6mm}
\par\noindent\begin{minipage}{\linewidth}
\textbf{14.} $\left[-10,4\right]$.\par
\end{minipage}\par\vspace{1.6mm}
\par\noindent\begin{minipage}{\linewidth}
\textbf{15.} a) $\left(5,+\infty\right)$ · b) $\left[3,+\infty\right)$ · c) $\left(-\infty,2\right)$ · d) $\left(-\infty,-3\right]$.\par
\end{minipage}\par\vspace{1.6mm}
\par\noindent\begin{minipage}{\linewidth}
\textbf{16.} a) $\left(8,+\infty\right)$ · b) $\left[-6,+\infty\right)$ · c) $\left[8,+\infty\right)$.\par
\end{minipage}\par\vspace{1.6mm}
\par\noindent\begin{minipage}{\linewidth}
\textbf{17.} $|T-21|\le1{,}5 \Rightarrow \left[19{,}5,\,22{,}5\right]$.\par
\end{minipage}\par\vspace{1.6mm}
\par\noindent\begin{minipage}{\linewidth}
\textbf{18.} a) $|\ell-50|\le0{,}2$ · b) $\left[49{,}8,\,50{,}2\right]$.\par
\end{minipage}\par\vspace{1.6mm}
\par\noindent\begin{minipage}{\linewidth}
\textbf{19.} a) $17\,\text{km}$ · b) $|p-3|\leq8 \Rightarrow \left[-5,11\right]$.\par
\end{minipage}\par\vspace{1.6mm}
\par\noindent\begin{minipage}{\linewidth}
\textbf{20.} a) $|x-7|>2$ · b) $\left(-\infty,5\right)\cup\left(9,+\infty\right)$ · c) No: $d\left(8{,}5,\,7\right)=1{,}5$, que no és més gran que $2$.\par
\end{minipage}\par\vspace{1.6mm}
\par\noindent\begin{minipage}{\linewidth}
\textbf{21.} $\left(1,10\right)$.\par
\end{minipage}\par\vspace{1.6mm}
\par\noindent\begin{minipage}{\linewidth}
\textbf{22.} $\left[-3,9\right]$.\par
\end{minipage}\par\vspace{1.6mm}
\par\noindent\begin{minipage}{\linewidth}
\textbf{23.} a) $A\cap B=\left(1,5\right)$, $A\cup B=\left[-2,8\right]$ · b) $A\cap B=\left(0,3\right]$, $A\cup B=\mathbb{R}$ · c) $A\cap B=\varnothing$, $A\cup B=\left[-4,-1\right]\cup\left[2,6\right]$ (no és un interval) · d) $A\cap B=\left[0,4\right)$, $A\cup B=\left(-3,+\infty\right)$.\par
\end{minipage}\par\vspace{1.6mm}
\par\noindent\begin{minipage}{\linewidth}
\textbf{24.} $|x-1|<4 \Rightarrow \left(-3,5\right)$; $|x+2|>1 \Rightarrow \left(-\infty,-3\right)\cup\left(-1,+\infty\right)$. La intersecció és $\left(-1,5\right)$.\par

\textbf{Bloc 3.}\par
\end{minipage}\par\vspace{1.6mm}
\par\noindent\begin{minipage}{\linewidth}
\textbf{25.} $2$ dec: truncament $2{,}71$, arrodoniment $2{,}72$. $4$ dec: truncament $2{,}7182$, arrodoniment $2{,}7183$.\par
\end{minipage}\par\vspace{1.6mm}
\par\noindent\begin{minipage}{\linewidth}
\textbf{26.} a) $2$ · b) $5$ · c) $5$ · d) $1$.\par
\end{minipage}\par\vspace{1.6mm}
\par\noindent\begin{minipage}{\linewidth}
\textbf{27.} $E_a=0{,}037$; $E_r\approx0{,}29\%$.\par
\end{minipage}\par\vspace{1.6mm}
\par\noindent\begin{minipage}{\linewidth}
\textbf{28.} a) $2\,\text{m}$ · b) $0{,}8\%$.\par
\end{minipage}\par\vspace{1.6mm}
\par\noindent\begin{minipage}{\linewidth}
\textbf{29.} Llapis: $6{,}67\%$; camp: $0{,}01\%$ — el camp és relativament més precís.\par
\end{minipage}\par\vspace{1.6mm}
\par\noindent\begin{minipage}{\linewidth}
\textbf{30.} $x\in\left[7{,}415,\;7{,}425\right)$. L'extrem inferior s'inclou, perquè $7{,}415$ s'arrodoneix a $7{,}42$ (la xifra següent és $5$ i fa pujar l'anterior); el superior no, perquè $7{,}425$ ja s'arrodoneix a $7{,}43$. L'error absolut màxim és, doncs, $0{,}005$, però només s'assoleix per sota.\par
\end{minipage}\par\vspace{1.6mm}
\par\noindent\begin{minipage}{\linewidth}
\textbf{31.} a) $E_a\leq0{,}005$, $E_r\leq\dfrac{0{,}005}{3{,}25}\approx0{,}15\%$ · b) $E_a\leq0{,}05$, $E_r\approx0{,}10\%$ · c) $E_a\leq5$, $E_r\approx0{,}42\%$ · d) $E_a\leq0{,}0005$, $E_r\approx1{,}32\%$. Fixa't que amb la mateixa precisió absoluta, com més petit és el valor mesurat, més gran és l'error relatiu.\par
\end{minipage}\par\vspace{1.6mm}
\par\noindent\begin{minipage}{\linewidth}
\textbf{32.} Calen $4$ decimals. En arrodonir a $n$ decimals, $E_a\leq0{,}5\cdot10^{-n}$, i per tant $E_r\leq\dfrac{0{,}5\cdot10^{-n}}{\pi}$. Amb $n=3$: $E_r\leq0{,}016\%$, que no arriba a la fita demanada; amb $n=4$: $E_r\leq0{,}0016\%<0{,}01\%$. Comprovació amb els valors reals: $\pi\approx3{,}142$ dona $E_r\approx0{,}013\%$ (insuficient) i $\pi\approx3{,}1416$ dona $E_r\approx0{,}0002\%$.\par
\end{minipage}\par\vspace{1.6mm}
\par\noindent\begin{minipage}{\linewidth}
\textbf{33.} $E_r\approx0{,}07\%$.\par

\textbf{Bloc 4.}\par
\end{minipage}\par\vspace{1.6mm}
\par\noindent\begin{minipage}{\linewidth}
\textbf{34.} a) $32$ · b) $81$ · c) $-81$ · d) $1$ · e) $0{,}01$ · f) $\tfrac49$.\par
\end{minipage}\par\vspace{1.6mm}
\par\noindent\begin{minipage}{\linewidth}
\textbf{35.} a) $3^6$ · b) $5^4$ · c) $2^6$ · d) $7^{-3}$ · e) $4^{-3}$.\par
\end{minipage}\par\vspace{1.6mm}
\par\noindent\begin{minipage}{\linewidth}
\textbf{36.} a) $729$ · b) $625$ · c) $64$ · d) $\tfrac1{343}$ · e) $\tfrac1{64}$.\par
\end{minipage}\par\vspace{1.6mm}
\par\noindent\begin{minipage}{\linewidth}
\textbf{37.} a) $36$ · b) $\tfrac{25}{9}$ · c) $16$ · d) $-8$.\par
\end{minipage}\par\vspace{1.6mm}
\par\noindent\begin{minipage}{\linewidth}
\textbf{38.} $2^0=1$.\par
\end{minipage}\par\vspace{1.6mm}
\par\noindent\begin{minipage}{\linewidth}
\textbf{39.} $3^0=1$.\par
\end{minipage}\par\vspace{1.6mm}
\par\noindent\begin{minipage}{\linewidth}
\textbf{40.} $\dfrac{2x}{3}$.\par
\end{minipage}\par\vspace{1.6mm}
\par\noindent\begin{minipage}{\linewidth}
\textbf{41.} $16$.\par
\end{minipage}\par\vspace{1.6mm}
\par\noindent\begin{minipage}{\linewidth}
\textbf{42.} a) $a^2$ · b) $x$ · c) $\dfrac{x^2}{y}$ · d) $2a$.\par
\end{minipage}\par\vspace{1.6mm}
\par\noindent\begin{minipage}{\linewidth}
\textbf{43.} a) $5{,}3\cdot10^4$ · b) $4{,}8\cdot10^{-4}$ · c) $7{,}2\cdot10^6$ · d) $6\cdot10^{-7}$.\par
\end{minipage}\par\vspace{1.6mm}
\par\noindent\begin{minipage}{\linewidth}
\textbf{44.} a) $320\,000$ · b) $0{,}009$ · c) $105$ · d) $0{,}0000067$.\par
\end{minipage}\par\vspace{1.6mm}
\par\noindent\begin{minipage}{\linewidth}
\textbf{45.} a) $6\cdot10^7$ · b) $3\cdot10^4$ · c) $4\cdot10^4$ · d) $3\cdot10^{-6}$.\par
\end{minipage}\par\vspace{1.6mm}
\par\noindent\begin{minipage}{\linewidth}
\textbf{46.} $8\cdot10^3$.\par
\end{minipage}\par\vspace{1.6mm}
\par\noindent\begin{minipage}{\linewidth}
\textbf{47.} $100\,000$ virus.\par
\end{minipage}\par\vspace{1.6mm}
\par\noindent\begin{minipage}{\linewidth}
\textbf{48.} $3\cdot10^5$ vegades.\par

\textbf{Bloc 5.}\par
\end{minipage}\par\vspace{1.6mm}
\par\noindent\begin{minipage}{\linewidth}
\textbf{49.} a) $5$ · b) $-3$ · c) no existeix · d) $2$ · e) $-2$ · f) $0$.\par
\end{minipage}\par\vspace{1.6mm}
\par\noindent\begin{minipage}{\linewidth}
\textbf{50.} a) $\pm7$ · b) $-2$ · c) cap · d) $\pm3$.\par
\end{minipage}\par\vspace{1.6mm}
\par\noindent\begin{minipage}{\linewidth}
\textbf{51.} a) $5^{1/3}$ · b) $2^{3/4}$ · c) $7^{1/2}$ · d) $x^{2/5}$.\par
\end{minipage}\par\vspace{1.6mm}
\par\noindent\begin{minipage}{\linewidth}
\textbf{52.} a) $\sqrt2$ · b) $\sqrt[3]{4}$ · c) $\sqrt[4]{x^3}$ · d) $\sqrt[4]{5}$.\par
\end{minipage}\par\vspace{1.6mm}
\par\noindent\begin{minipage}{\linewidth}
\textbf{53.} a) $\sqrt[3]{2}$ · b) $\sqrt3$ · c) $\sqrt2$ · d) $\sqrt3$.\par
\end{minipage}\par\vspace{1.6mm}
\par\noindent\begin{minipage}{\linewidth}
\textbf{54.} a) $5\sqrt3$ · b) $3\sqrt[3]{2}$ · c) $10\sqrt2$ · d) $4\sqrt[3]{2}$ · e) $x^2\sqrt x$ · f) $2x\sqrt[3]{x}$ · g) $\dfrac{x^2}{2}$ · h) $3a\sqrt{2a}$.\par
\end{minipage}\par\vspace{1.6mm}
\par\noindent\begin{minipage}{\linewidth}
\textbf{55.} a) $\sqrt{18}$ · b) $\sqrt[3]{40}$ · c) $\sqrt{x^3}$ · d) $\sqrt[3]{128}$ · e) $\sqrt{a}$ · f) $\sqrt5$.\par
\end{minipage}\par\vspace{1.6mm}
\par\noindent\begin{minipage}{\linewidth}
\textbf{56.} a) $\dfrac{\sqrt{13}}{6}$ · b) $\dfrac{3}{2}\sqrt{x}$ · c) $3\sqrt[3]{x^2+1}$ · d) $|2x+1|$ (atenció: \textbf{no} és $2x+1$, perquè $x$ pot ser menor que $-\tfrac12$).\par
\end{minipage}\par\vspace{1.6mm}
\par\noindent\begin{minipage}{\linewidth}
\textbf{57.} $\sqrt[6]{5}<\sqrt2<\sqrt[3]{3}<\sqrt[3]{4}$ (índex comú $6$: $\sqrt[6]{5}$, $\sqrt[6]{8}$, $\sqrt[6]{9}$, $\sqrt[6]{16}$).\par
\end{minipage}\par\vspace{1.6mm}
\par\noindent\begin{minipage}{\linewidth}
\textbf{58.} a) $6$ · b) $2$ · c) $10$ · d) $2$.\par
\end{minipage}\par\vspace{1.6mm}
\par\noindent\begin{minipage}{\linewidth}
\textbf{59.} a) $5$ · b) $3$ · c) $3$ · d) $2$.\par
\end{minipage}\par\vspace{1.6mm}
\par\noindent\begin{minipage}{\linewidth}
\textbf{60.} a) $\sqrt[6]{32}$ · b) $\sqrt[6]{3}$.\par
\end{minipage}\par\vspace{1.6mm}
\par\noindent\begin{minipage}{\linewidth}
\textbf{61.} a) $2$ · b) $\sqrt[12]{x}$ · c) $\sqrt[3]{25}$ · d) $49$.\par
\end{minipage}\par\vspace{1.6mm}
\par\noindent\begin{minipage}{\linewidth}
\textbf{62.} a) $\sqrt[4]{125}$ · b) $\sqrt2$ · c) $\sqrt[6]{2}$ · d) $\sqrt[3]{9}$.\par
\end{minipage}\par\vspace{1.6mm}
\par\noindent\begin{minipage}{\linewidth}
\textbf{63.} a) $8\sqrt2$ · b) $5\sqrt[3]{3}$ · c) $-3\sqrt5$.\par
\end{minipage}\par\vspace{1.6mm}
\par\noindent\begin{minipage}{\linewidth}
\textbf{64.} a) $5\sqrt2$ · b) $\sqrt2$ · c) $13\sqrt3$.\par
\end{minipage}\par\vspace{1.6mm}
\par\noindent\begin{minipage}{\linewidth}
\textbf{65.} a) $4\sqrt5$ · b) $5\sqrt[3]{a}$ · c) $6\sqrt x$.\par
\end{minipage}\par\vspace{1.6mm}
\par\noindent\begin{minipage}{\linewidth}
\textbf{66.} a) $\dfrac{5\sqrt3}{3}$ · b) $\dfrac{2\sqrt5}{5}$ · c) $3\sqrt2$ · d) $\dfrac{\sqrt7}{7}$.\par
\end{minipage}\par\vspace{1.6mm}
\par\noindent\begin{minipage}{\linewidth}
\textbf{67.} a) $\dfrac{2\sqrt3}{3}$ · b) $\dfrac{\sqrt3}{2}$ · c) $\dfrac{2\sqrt5}{3}$.\par
\end{minipage}\par\vspace{1.6mm}
\par\noindent\begin{minipage}{\linewidth}
\textbf{68.} a) $\sqrt5+2$ · b) $\dfrac{\sqrt7-1}{2}$ · c) $\dfrac{6+2\sqrt2}{7}$.\par
\end{minipage}\par\vspace{1.6mm}
\par\noindent\begin{minipage}{\linewidth}
\textbf{69.} a) $5\sqrt3+5\sqrt2$ · b) $\dfrac{\sqrt n}{n}$ · c) $\sqrt a-\sqrt b$.\par
\end{minipage}\par\vspace{1.6mm}
\par\noindent\begin{minipage}{\linewidth}
\textbf{70.} $0$.\par
\end{minipage}\par\vspace{1.6mm}
\par\noindent\begin{minipage}{\linewidth}
\textbf{71.} $c=5\sqrt2$.\par
\end{minipage}\par\vspace{1.6mm}
\par\noindent\begin{minipage}{\linewidth}
\textbf{72.} a) $18\,\text{cm}^2$ · b) $6\,\text{cm}$.\par
\end{minipage}\par\vspace{1.6mm}
\par\noindent\begin{minipage}{\linewidth}
\textbf{73.} $\sqrt{13}\approx3{,}606\neq5$.\par
\end{minipage}\par\vspace{1.6mm}
\par\noindent\begin{minipage}{\linewidth}
\textbf{74.} $x^2$.\par

\textbf{Bloc 6.}\par
\end{minipage}\par\vspace{1.6mm}
\par\noindent\begin{minipage}{\linewidth}
\textbf{75.} a) $\log_2 32=5$ · b) $\log 1000=3$ · c) $\log_5 1=0$ · d) $\log_3\tfrac19=-2$.\par
\end{minipage}\par\vspace{1.6mm}
\par\noindent\begin{minipage}{\linewidth}
\textbf{76.} a) $2^4=16$ · b) $10^2=100$ · c) $3^0=1$ · d) $5^1=5$.\par
\end{minipage}\par\vspace{1.6mm}
\par\noindent\begin{minipage}{\linewidth}
\textbf{77.} a) $3$ · b) $4$ · c) $4$ · d) $0$ · e) $1$ · f) $1$.\par
\end{minipage}\par\vspace{1.6mm}
\par\noindent\begin{minipage}{\linewidth}
\textbf{78.} a) $-2$ · b) $-3$ · c) $\tfrac12$ · d) $\tfrac23$.\par
\end{minipage}\par\vspace{1.6mm}
\par\noindent\begin{minipage}{\linewidth}
\textbf{79.} a) $2+3=5$ · b) $\log5+\log20=2$ · c) $\log_3x+\log_3y$.\par
\end{minipage}\par\vspace{1.6mm}
\par\noindent\begin{minipage}{\linewidth}
\textbf{80.} a) $5-2=3$ · b) $3-1=2$ · c) $\log_5x-2$.\par
\end{minipage}\par\vspace{1.6mm}
\par\noindent\begin{minipage}{\linewidth}
\textbf{81.} a) $15$ · b) $-4$ · c) $7$ · d) $n$.\par
\end{minipage}\par\vspace{1.6mm}
\par\noindent\begin{minipage}{\linewidth}
\textbf{82.} a) $1$ · b) $1$ · c) $1$.\par
\end{minipage}\par\vspace{1.6mm}
\par\noindent\begin{minipage}{\linewidth}
\textbf{83.} $1{,}5$.\par
\end{minipage}\par\vspace{1.6mm}
\par\noindent\begin{minipage}{\linewidth}
\textbf{84.} $-0{,}5$.\par
\end{minipage}\par\vspace{1.6mm}
\par\noindent\begin{minipage}{\linewidth}
\textbf{85.} a) $\log_2 15$ · b) $\log 5$ · c) $\log_5 8$ · d) $\log_3 3=1$.\par
\end{minipage}\par\vspace{1.6mm}
\par\noindent\begin{minipage}{\linewidth}
\textbf{86.} $\log_2 4=2$.\par
\end{minipage}\par\vspace{1.6mm}
\par\noindent\begin{minipage}{\linewidth}
\textbf{87.} a) $\dfrac{\log 7}{\log 2}$ · b) $\dfrac{\log 12}{\log 5}$.\par
\end{minipage}\par\vspace{1.6mm}
\par\noindent\begin{minipage}{\linewidth}
\textbf{88.} a) $3\log_a x+2\log_a y$ · b) $2\log_a x-3\log_a y$ · c) $\log_a x+\log_a y-\log_a z$ · d) $\tfrac12\log_a x$.\par
\end{minipage}\par\vspace{1.6mm}
\par\noindent\begin{minipage}{\linewidth}
\textbf{89.} a) $3\log_a x+\tfrac12\log_a y-2\log_a z$ · b) $\tfrac12\left(\log_a x-\log_a y\right)$ · c) $2\log_a x+\tfrac13\log_a y-\log_a z$ · d) $\log_a x+\tfrac12\log_a y+\tfrac12\log_a z$.\par
\end{minipage}\par\vspace{1.6mm}
\par\noindent\begin{minipage}{\linewidth}
\textbf{90.} a) $\log_a\left(x^2y^3\right)$ · b) $\log_a\dfrac{x}{\sqrt y}$ · c) $\log_a\dfrac{xy}{z}$ · d) $\log_a x^4$.\par
\end{minipage}\par\vspace{1.6mm}
\par\noindent\begin{minipage}{\linewidth}
\textbf{91.} a) $\log_a\dfrac{x^3}{yz^2}$ · b) $\log_a\sqrt{xy}$ · c) $\log_a\dfrac{x^2}{\sqrt[3]{y}}$ · d) $\log_a\dfrac{\sqrt[3]{xy^2}}{z}$.\par
\end{minipage}\par\vspace{1.6mm}
\par\noindent\begin{minipage}{\linewidth}
\textbf{92.} a) $m+n$ · b) $2m+n$ · c) $n-m$ · d) $\tfrac{m}{2}$ · e) $m+2n$ · f) $3m-2n$.\par
\end{minipage}\par\vspace{1.6mm}
\par\noindent\begin{minipage}{\linewidth}
\textbf{93.} a) $0{,}699$ · b) $0{,}602$ · c) $1{,}699$ · d) $-0{,}699$ · e) $0{,}1505$ · f) $1{,}602$.\par
\end{minipage}\par\vspace{1.6mm}
\par\noindent\begin{minipage}{\linewidth}
\textbf{94.} $\tfrac12 L$ (ja que $2L-\tfrac12L-L=\tfrac12L$).\par
\end{minipage}\par\vspace{1.6mm}
\par\noindent\begin{minipage}{\linewidth}
\textbf{95.} Pel canvi de base, $\log_b a=\dfrac{\log_a a}{\log_a b}=\dfrac{1}{\log_a b}$. Per tant $\log_a b\cdot\log_b a=\log_a b\cdot\dfrac{1}{\log_a b}=1$.\par
\end{minipage}\par\vspace{1.6mm}
\par\noindent\begin{minipage}{\linewidth}
\textbf{96.} $\log_{a^2}x=\dfrac{\log_a x}{\log_a a^2}=\dfrac{\log_a x}{2}=\tfrac12\log_a x$.\par
\end{minipage}\par\vspace{1.6mm}
\par\noindent\begin{minipage}{\linewidth}
\textbf{97.} $\approx8{,}6$ anys.\par
\end{minipage}\par\vspace{1.6mm}
\par\noindent\begin{minipage}{\linewidth}
\textbf{98.} $M-M_1=2$.\par

\textbf{Bloc 7.}\par
\end{minipage}\par\vspace{1.6mm}
\par\noindent\begin{minipage}{\linewidth}
\textbf{99.} $4$.\par
\end{minipage}\par\vspace{1.6mm}
\par\noindent\begin{minipage}{\linewidth}
\textbf{100.} $2\cdot10^5$.\par
\end{minipage}\par\vspace{1.6mm}
\par\noindent\begin{minipage}{\linewidth}
\textbf{101.} $\tfrac1{15}$.\par
\end{minipage}\par\vspace{1.6mm}
\par\noindent\begin{minipage}{\linewidth}
\textbf{102.} $\left(-1,7\right)$.\par
\end{minipage}\par\vspace{1.6mm}
\par\noindent\begin{minipage}{\linewidth}
\textbf{103.} $\approx0{,}015\%$.\par
\end{minipage}\par\vspace{1.6mm}
\par\noindent\begin{minipage}{\linewidth}
\textbf{104.} a) $16\,000$ · b) $7$ hores.\par
\end{minipage}\par\vspace{1.6mm}
\par\noindent\begin{minipage}{\linewidth}
\textbf{105.} a) $5\,\text{m}$ · b) $5\sqrt[3]{2}\,\text{m}$.\par
\end{minipage}\par\vspace{1.6mm}
\par\noindent\begin{minipage}{\linewidth}
\textbf{106.} $80\,\text{dB}$.\par
\end{minipage}\par\vspace{1.6mm}
\par\noindent\begin{minipage}{\linewidth}
\textbf{107.} $4$.\par
\end{minipage}\par\vspace{1.6mm}
\par\noindent\begin{minipage}{\linewidth}
\textbf{108.} a) $A=6$, $B=3$ · b) $\left(3,9\right)$.\par
\end{minipage}\par\vspace{1.6mm}
\par\noindent\begin{minipage}{\linewidth}
\textbf{109.} a) F ($a=-2$: $\sqrt{4}=2\neq-2$; en realitat $\sqrt{a^2}=|a|$) · b) F ($a=1$, $b=-1$: $0\neq2$) · c) F ($x=y=1$: $\log2\neq0$) · d) C · e) C · f) F ($a=b=1$: $4\neq2$).\par
\end{minipage}\par\vspace{1.6mm}
\par\noindent\begin{minipage}{\linewidth}
\textbf{110.} Si $T_2=2T_1$, aleshores $2\pi\sqrt{\tfrac{L_2}{g}}=2\cdot2\pi\sqrt{\tfrac{L_1}{g}}$. Simplificant $2\pi$ i $g$: $\sqrt{L_2}=2\sqrt{L_1}$. Elevant al quadrat, $L_2=4L_1$.\par
\end{minipage}\par\vspace{1.6mm}
\end{multicols}
\end{document}
